% Pursuing Stacks (cotes 134-2 to 134-8): archive cote 134-6, site page 14 (PDF page 15), typescript page 333.
% typescript. Draft compared: ps5.tex:2041-2123.
% AI-transcribed by Claude Opus 5.5 (claude-opus-5-5, adaptive thinking, effort high), Message Batches API, batch ps3 (msgbatch_01XdFmuVVBAcgc7EFxdtCqsj), 2026-11-12; not yet checked by eye.
\scanpage{134-6/14}
\mspage{333}

When A is a finite product of copies taken from among the three
standard test categories \hand{$\Delta$},\hand{$\square$} ,\hand{$\bigcirc$} ,the standard chain complex structure
on multicomplexes does furnish us with a \uline{``canonical''} \uline{abelianizator}
for A, which we may denote by $L_{\cdot}^{A}$ (as we did yesterday for A =\hand{$\Delta$} ).
This ``standard'' abelianizator has some very remarkable extra features
which I would like to pin down, \xout{and} which had caused our rather narrow
focus in the notes of two days ago (section 91).

a) There is a ``dimension map''
\[ (7) \qquad \dim:\ \mathrm{Ob}\,A \longrightarrow \underline{N} \quad . \]
It can be described (in the particular case above at any rate) in terms
of the intrinsic category structure of A, by associating to every a in
A the ordered set
\[ (8) \qquad i(a) \;=\; \text{set of subobjects of a in A} \hand{\unsure{.}} \]
(NB not to be confused with subobjects of a in A\^{}, namely sieves in
a). \deleted{which}\inserted{This} is an ordered set with \inserted{a} largest object \xout{a} (namely a itself),
and which turns out to be finite (in the particular case\xout{s} considered),
hence of finite combinatorial dimension (equal to the dimension of the
geometrical realization $|i(a)|$ ), and we have
\[ (9) \qquad \dim\,(a) \;=\; \dim\ i(a) \quad . \]

b) The n.th component $L_n$ of $L_{\cdot}^{A} = L.$ is given by
\[ (10) \qquad L_n \;=\; \hand{$\displaystyle\bigoplus$}_{\dim(a)=n} a \quad , \]
(where the direct sum of course is taken in Add($A^{o}$) as in (2) above\hand{,})
which makes a sense when we assume (as is the case in our example)
that the map (7) \inserted[typed]{is ``finite'', i.e.\hand{,}} \deleted{\xout{\illegible{3}}} has finite fibers. For instance, in
all \xout{\illegible{1}} our ``standard'' cases, there is just \uline{one} object of A which is
of dimension 0, \xout{namely} and this is also the final object.

c) The differential operator
\[ \hand{(11)} \qquad d_n:\ L_n \longrightarrow L_{n-1} \]
can be obtained in the following way. We have \xout{t}only to describe $d_n$ on
each summand a of $L_n$ i.e. by (10) on each a \hand{in} \deleted{\xout{\illegible{4}}}
\[ \xout{\illegible{1}} \ (12) \qquad A_n = \hand{$\{$}\, a \hand{$\in$} \mathrm{Ob}\,A \hand{$|$} \dim a = n \hand{$\}$}. \]
In view of (3), this restriction $d_n$ \hand{$|$} a can be described as a linear
combination of elements in the disjoint sum of the sets
\[ \mathrm{Hom}(b,a) \quad , \quad \text{with}\ \ b \hand{$\in$} An\text{-}1 \quad \hand{(a $\in A_n$ fixed).} \]

\margin{14}
